\documentclass[aspectratio=119,12pt]{beamer}

\usepackage{amsmath}
\usepackage{tikz}
\usepackage{pgfplots}
\usepackage{pgfpages}
% \usepackage{svg}

\pgfplotsset{compat=1.18}

\usetheme{default}

\setbeamertemplate{navigation symbols}{}

\pgfpagesuselayout{8 on 1}[
    a4paper,
    border shrink=2mm
]

\setbeamertemplate{background}{
    \begin{tikzpicture}[remember picture,overlay]
        \draw[gray, line width=0.5pt]
            ([xshift=3pt,yshift=-3pt]current page.north west)
            rectangle
            ([xshift=-3pt,yshift=3pt]current page.south east);
    \end{tikzpicture}
}

\begin{document}

%------------------------------------------------
\begin{frame}
    \centering
     \textbf{Tallmengder}

     \begin{itemize}
         \item $ \mathbb{N} $ være naturlige tall 
             \begin{itemize}
                 \item $1,2,3,4,5 \ldots$
                 \item $\mathbb{N}_0 = 0,1,2,3,4 \ldots $
             \end{itemize}
         \item $ \mathbb{Z} $ alle hele negative og positive tall 
             \begin{itemize}
                 \item $-3,-2,-1,0,1,2,3,4$
             \end{itemize}
                 
         \item $ \mathbb{Q}  $ rasjonelle tall: Alle brøker 
             \begin{itemize}
                 \item  $\frac{1}{3} , \frac{1}{5}, \frac{10}{150} , \ldots$ 
             \end{itemize}

         \item $ \mathbb{R} $ være relle tall. 
             \begin{itemize}
                 \item Alle tall på tall-linjen. 
                 \item $1.01, \frac{1}{2} , \pi, 99,99 $ . 
             \end{itemize}
     \end{itemize}




\end{frame}

%------------------------------------------------
\begin{frame}{Vertikal tavle: Oppgave 1}

    La $a,b,c$ og $d$ være tilfeldige tall. Vis at \begin{equation}
 \frac{a}{b} :  \frac{c}{d} = \frac{a}{b} \cdot \frac{d}{c}         
    \end{equation}  

\end{frame}



\begin{frame}{Vertikal tavle: Oppgave 2}

    La $a,b,c$ og $d$ være tilfeldige tall.

    Vis at  \begin{equation}
        (a +b)^{2} = a^2 + b^2 + 2ab
    \end{equation}

    Bruk forrige resultat til å vise at
    \begin{equation}
        (a +b)(a-b) = a^2 - b^2 
    \end{equation}


\end{frame}


\begin{frame}{Vertikal tavle: Oppgave 3}


    La $x$ være et ukjent tall. 
    Prøv å skrive \begin{equation}
        x^2 - 4    
    \end{equation}
    i "parantesform". Hint: $(a+b)(a-b) = a^2 - b^2$.  


\vskip 0.4cm
    \begin{itemize}
        \item Husker dere faktorisering? $18 = 3\cdot 3\cdot 2$? Det kalles å skrive et utrykk som produkt (altså ganging) av faktorer (feks et tall eller noe i en parantes).   
    \end{itemize}




\end{frame}



\begin{frame}{Sykepleierlønn og boligpriser}

    \begin{columns}[T]

        % Sykepleierlønn
        \begin{column}{0.48\textwidth}
            \centering
            \textbf{Sykepleierlønn i Norge}

            \vspace{0.4cm}

            \begin{tabular}{lr}
                \hline
                \textbf{År} & \textbf{Kr/mnd} \\
                \hline
                2020 & 47\,020 \\
                2021 & 49\,440 \\
                2022 & 51\,020 \\
                2023 & 53\,880 \\
                2024 & 56\,710 \\
                2025 & 58\,290 \\
                \hline
            \end{tabular}
        \end{column}

        % Boligpriser
        \begin{column}{0.48\textwidth}
            \centering
            \textbf{Eneboliger i Vestvågøy}

            \vspace{0.4cm}

            \begin{tabular}{lr}
                \hline
                \textbf{År} & \textbf{Kr/m$^2$} \\
                \hline
                2020 & 20\,705 \\
                2021 & 22\,460 \\
                2022 & 21\,823 \\
                2023 & 23\,483 \\
                2024 & 25\,301 \\
                \hline
            \end{tabular}
        \end{column}

    \end{columns}

    \vspace{0.6cm}

    \begin{center}
        \footnotesize Kilde: Statistisk sentralbyrå (SSB)
    \end{center}

\end{frame}


\end{document}
